% !TEX root = task-1-3-1.tex
% Editable task bank. Append new BankTask blocks at the end of the document.
\documentclass[12pt,a4paper]{article}
\usepackage{physics-textbook}
\usepackage{graphicx,needspace}
\usetikzlibrary{decorations.pathmorphing}
\newif\ifShowSolutions
\ifdefined\TutorEdition
  \ShowSolutionstrue
\else
  \ShowSolutionsfalse % Change to \ShowSolutionstrue to display solutions.
\fi
\newcounter{banktask}
\newsavebox{\BankTaskBox}
\newcommand{\BankTaskContent}[3]{%
  \noindent\textbf{\thebanktask.}\quad #2\par
  \ifShowSolutions
    \Needspace{50mm}% Keep the solution heading with its opening text or diagram.
    #3\par
    \ifdefined\TutorEdition
      \ifcsname TutorSource#1\endcsname
        \nopagebreak\smallskip{\small\csname TutorSource#1\endcsname\par}%
      \fi
    \fi
  \fi
}
\newcommand{\BankTask}[3]{%
  \refstepcounter{banktask}\par\medskip
  \sbox{\BankTaskBox}{\begin{minipage}{\linewidth}%
    \BankTaskContent{#1}{#2}{#3}\end{minipage}}%
  \ifdim\dimexpr\ht\BankTaskBox+\dp\BankTaskBox\relax>\dimexpr\textheight-20mm\relax
    \Needspace{35mm}\BankTaskContent{#1}{#2}{#3}%
  \else
    \Needspace{\dimexpr\ht\BankTaskBox+\dp\BankTaskBox\relax}%
    \noindent\usebox{\BankTaskBox}\par
  \fi
  \medskip
}
\begin{document}
\BookSubsection{1.3.1}{Зависимость скорости от времени}


% Source 081272: https://ege.fipi.ru/bank/questions.php?search=1&proj=BA1F39653304A5B041B656915DC36B38&qid=081272
\BankTask{081272}{
На рисунке приведены графики зависимости проекции скорости от времени для легкового автомобиля ($\mathrm{I}$) и микроавтобуса ($\mathrm{II}$), движущихся по прямой дороге, вдоль которой и направлена ось $Ox$. Определите отношение модулей ускорений $\dfrac{a_{\mathrm I}}{a_{\mathrm{II}}}$.
\begin{center}
\begin{tikzpicture}[x=1.4cm,y=0.14cm,>=stealth]
\draw[black!85,thin,xstep=0.5,ystep=5] (0,0) grid (4.5,25);
\draw[->] (0,0)--(5.15,0) node[right] {$t,\,\text{с}$};
\draw[->] (0,0)--(0,29) node[above] {$\upsilon_x,\,\text{м}/\text{с}$};
\node[below left] at (0,0) {$0$};
\node[left] at (0,20) {$20$};
\foreach \t in {2,4} \node[below] at (\t,0) {$\t$};
\draw[very thick] (0,20)--(4,0);
\draw[very thick] (0,0)--(4.5,11.25);
\node at (1.25,17.3) {$\mathrm I$};
\node at (3.8,13.3) {$\mathrm{II}$};
\end{tikzpicture}
\end{center}
}{
\smallskip\noindent\textbf{Решение от GPT}\par
Модуль ускорения равен модулю наклона графика $\upsilon_x(t)$. За $4\,\text{с}$ скорость автомобиля изменяется от $20$ до $0\,\text{м}/\text{с}$, а микроавтобуса --- от $0$ до $10\,\text{м}/\text{с}$. Тогда $a_{\mathrm I}=20/4=5\,\text{м}/\text{с}^2$, $a_{\mathrm{II}}=10/4=2{,}5\,\text{м}/\text{с}^2$, и $\dfrac{a_{\mathrm I}}{a_{\mathrm{II}}}=2$.
}

% Source 277542: https://ege.fipi.ru/bank/questions.php?search=1&proj=BA1F39653304A5B041B656915DC36B38&qid=277542
\BankTask{277542}{
На рисунке показан график зависимости проекции $\upsilon_x$ скорости тела от времени $t$. Какова проекция $a_x$ ускорения этого тела в момент времени $2\,\text{с}$? Результат выразите в $\text{м}/\text{с}^2$.
\begin{center}
\begin{tikzpicture}[x=1.3cm,y=0.23cm,>=stealth]
\draw[thin,black!65,xstep=0.5,ystep=2] (0,-6) grid (3,6);
\draw[->] (0,0)--(4.0,0) node[right] {$t,\,\text{с}$};
\draw[->] (0,-6)--(0,8) node[above] {$\upsilon_x,\,\text{м}/\text{с}$};
\node[below] at (1,0) {$1$};
\node[below] at (2,0) {$2$};
\node[below] at (3,0) {$3$};
\node[left] at (0,-6) {$-6$};
\node[left] at (0,-4) {$-4$};
\node[left] at (0,-2) {$-2$};
\node[left] at (0,0) {$0$};
\node[left] at (0,2) {$2$};
\node[left] at (0,4) {$4$};
\node[left] at (0,6) {$6$};
\draw[very thick] (0,4)--(0.5,0)--(1,6)--(1.5,2)--(2.5,2)--(3,-6);
\end{tikzpicture}
\end{center}
}{
\smallskip\noindent\textbf{Решение от GPT}\par
При $t=2\,\text{с}$ график горизонтален: $a_x=\Delta v_x/\Delta t=0\,\text{м}/\text{с}^2$.
}

% Source eDB1FA: https://ege.fipi.ru/bank/questions.php?search=1&proj=BA1F39653304A5B041B656915DC36B38&qid=eDB1FA
\BankTask{eDB1FA}{
На рисунке приведён график зависимости проекции $\upsilon_x$ скорости тела
от времени $t$.

Определите проекцию $a_{x}$ ускорения этого тела в интервале времени
от $0$ до $2\,\text{с}$. Учитывайте знак проекции. Результат выразите в $\text{м}/\text{с}^2$.
\begin{center}
\begin{tikzpicture}[x=0.48cm,y=0.18cm,>=stealth]
\draw[thin,black!65,xstep=1,ystep=3] (0,-9) grid (9,9);
\draw[->] (0,0)--(11,0) node[right] {$t,\,\text{с}$};
\draw[->] (0,-9)--(0,12) node[above] {$\upsilon_x,\,\text{м}/\text{с}$};
\node[below] at (1,0) {$1$};
\node[below] at (2,0) {$2$};
\node[below] at (3,0) {$3$};
\node[below] at (4,0) {$4$};
\node[below] at (5,0) {$5$};
\node[below] at (6,0) {$6$};
\node[below] at (7,0) {$7$};
\node[below] at (8,0) {$8$};
\node[below] at (9,0) {$9$};
\node[left] at (0,-9) {$-9$};
\node[left] at (0,-6) {$-6$};
\node[left] at (0,-3) {$-3$};
\node[left] at (0,0) {$0$};
\node[left] at (0,3) {$3$};
\node[left] at (0,6) {$6$};
\node[left] at (0,9) {$9$};
\draw[very thick] (0,-9)--(3,4.5)--(5.5,4.5)--(9,-6);
\end{tikzpicture}
\end{center}
}{
\smallskip\noindent\textbf{Решение от GPT}\par
На первых $2\,\text{с}$ $v_x$ изменяется от $-9$ до $0\,\text{м}/\text{с}$. Поэтому $a_x=(0-(-9))/2=4{,}5\,\text{м}/\text{с}^2$.
}

% Source B61411: https://ege.fipi.ru/bank/questions.php?search=1&proj=BA1F39653304A5B041B656915DC36B38&qid=B61411
\BankTask{B61411}{
На рисунке показан график зависимости проекции $\upsilon_x$ скорости тела от времени $t$. Какова проекция $a_x$ ускорения этого тела в интервале времени от $2$ до $3\,\text{с}$? Учитывайте знак проекции. Результат выразите в $\text{м}/\text{с}^2$.
\begin{center}
\begin{tikzpicture}[x=0.7cm,y=0.175cm,>=stealth]
\draw[thin,black!65,xstep=1,ystep=4] (0,-12) grid (6,12);
\draw[->] (0,0)--(8,0) node[right] {$t,\,\text{с}$};
\draw[->] (0,-13.5)--(0,17.5) node[above] {$\upsilon_x,\,\text{м}/\text{с}$};
\node[below] at (2,0) {$2$};
\node[below] at (4,0) {$4$};
\node[below] at (6,0) {$6$};
\node[left] at (0,-12) {$-12$};
\node[left] at (0,-8) {$-8$};
\node[left] at (0,-4) {$-4$};
\node[left] at (0,0) {$0$};
\node[left] at (0,4) {$4$};
\node[left] at (0,8) {$8$};
\node[left] at (0,12) {$12$};
\draw[very thick] (0,12)--(2,4)--(3,-12)--(4,-12)--(5,0)--(6,-8);
\end{tikzpicture}
\end{center}
}{
\smallskip\noindent\textbf{Решение от GPT}\par
$a_x=(v_x(3)-v_x(2))/(3-2)=(-12-4)/1=-16\,\text{м}/\text{с}^2$.
}

% Source eF04e1: https://ege.fipi.ru/bank/questions.php?search=1&proj=BA1F39653304A5B041B656915DC36B38&qid=eF04e1
\BankTask{eF04e1}{
На рисунке показан график зависимости проекции $\upsilon_x$ скорости тела от времени $t$. Какова проекция $a_x$ ускорения этого тела в интервале времени от $0$ до $0{,}5\,\text{с}$? Учитывайте знак проекции. Результат выразите в $\text{м}/\text{с}^2$.
\begin{center}
\begin{tikzpicture}[x=1cm,y=0.25cm,>=stealth]
\draw[thin,black!65,xstep=0.5,ystep=2] (0,-6) grid (3,6);
\draw[->] (0,0)--(4.0,0) node[right] {$t,\,\text{с}$};
\draw[->] (0,-6)--(0,8) node[above] {$\upsilon_x,\,\text{м}/\text{с}$};

\node[below] at (1,0) {$1$};
\node[below] at (2,0) {$2$};
\node[below] at (3,0) {$3$};
\node[left] at (0,-6) {$-6$};
\node[left] at (0,-4) {$-4$};
\node[left] at (0,-2) {$-2$};
\node[left] at (0,0) {$0$};
\node[left] at (0,2) {$2$};
\node[left] at (0,4) {$4$};
\node[left] at (0,6) {$6$};
\draw[very thick] (0,-6)--(0.5,2)--(1.5,2)--(2,6)--(2.5,0)--(3,4);
\end{tikzpicture}
\end{center}
}{
\smallskip\noindent\textbf{Решение от GPT}\par
$a_x=(2-(-6))/0{,}5=16\,\text{м}/\text{с}^2$.
}

% Source D82e7B: https://ege.fipi.ru/bank/questions.php?search=1&proj=BA1F39653304A5B041B656915DC36B38&qid=D82e7B
\BankTask{D82e7B}{
Два тела, каждое массой $2\,\text{кг}$, движутся на одной прямой, вдоль которой направлена ось $O x$ На рисунке приведены графики зависимости проекций $\upsilon_x$ их скорости от времени $t.$
\begin{center}\begin{tikzpicture}[x=.225cm,y=.45cm,>=stealth]
\draw[thin,dashed,xstep=2,ystep=1] (0,0) grid (20,8);
\draw[->] (0,0)--(22,0) node[right] {$t,\,\text{с}$};
\draw[->] (0,0)--(0,9) node[above] {$\upsilon_x,\,\text{м}/\text{с}$};
\foreach \x in {0,4,8,12,16} \node[below] at (\x,0) {$\x$};
\foreach \y in {2,4,6} \node[left] at (0,\y) {$\y$};
\draw[very thick] (0,7)--(20,{19/7}); \node at (6,6.8) {$1$};
\draw[very thick] (0,.5)--(20,5.5); \node at (8,1.7) {$2$};
\end{tikzpicture}\end{center}
\begin{enumerate}
\item Равнодействующая сил, действующих на тело $1$, остаётся постоянной в течение всего времени наблюдения.
\item В промежутке времени от $0$ до $14\,\text{с}$ тела двигались навстречу друг другу.
\item Путь, пройденный телом $1$ за промежуток времени от $0$ до $18\,\text{с}$, больше пути, пройденного телом $2$ за тот же промежуток времени.
\item Кинетическая энергия тела $2$ за промежуток времени от $6$ до $14\,\text{с}$ увеличилась в $2$ раза.
\item Модуль импульса тела $2$ в момент времени $t=10\,\text{с}$ равен $10\,\text{кг}\cdot\text{м}/\text{с}$.
\end{enumerate}
}{
\smallskip\noindent\textbf{Решение от GPT}\par
Линейные графики задают $v_{1x}=7-\frac{3}{14}t$ и $v_{2x}=0{,}5+0{,}25t$ в СИ. Равнодействующая на тело $1$ постоянна; обе скорости положительны. За первые $18\,\text{с}$ пути $s_1=7\cdot18-\frac{3}{28}\cdot18^2=\frac{639}{7}\,\text{м}\approx91{,}3\,\text{м}$ и $s_2=(0{,}5+5)18/2=49{,}5\,\text{м}$, значит $s_1>s_2$. Для тела $2$ $E_{\text{к}}(14)/E_{\text{к}}(6)=(4/2)^2=4$, $p(10)=2\cdot3=6\,\text{кг}\cdot\text{м}/\text{с}$. Ответ: $13$.
}

\end{document}
